\documentclass[11pt]{article}

\usepackage[margin=1in]{geometry}
\usepackage{hyperref}
\usepackage{xcolor}
\usepackage{listings}
\usepackage{booktabs}
\usepackage{parskip}

\lstset{
  basicstyle=\ttfamily\small,
  breaklines=true,
  frame=single,
  backgroundcolor=\color{gray!8},
  columns=fullflexible,
}

\hypersetup{colorlinks=true, linkcolor=blue, urlcolor=blue}

\title{A Field Guide to the DeepScratch Playground UI\\
\large For people who think in tensors, not pixels}
\author{}
\date{\today}

\begin{document}
\maketitle
\tableofcontents

\section{Why this document exists}

You know exactly what an encoder, a head, a loss function, and a train/val/test
split are. You have never had a reason to know what a ``DOM'' is or why a
browser cares about something called an ``event listener.'' This guide exists
so you can look at the playground, know something is wrong or something needs
to change, and go directly to the right file and the right ten lines --
without learning web development.

The playground is small on purpose: two files you will actually edit, one
file you will almost never touch, and one Python backend that is mostly
plumbing you already understand (it reads a CSV, writes a YAML file, and
calls \lstinline{subprocess.Popen} on \lstinline{entrypoint.py}).

\section{The Mental Model}

If you squint, the frontend is a training loop wearing a costume:

\begin{center}
\begin{tabular}{@{}ll@{}}
\toprule
\textbf{Training loop concept} & \textbf{Playground equivalent} \\
\midrule
Model weights (the state you mutate) & the \lstinline{state} object in \lstinline{app.js} \\
Forward pass (state $\to$ output) & the \lstinline{render*()} functions (state $\to$ DOM) \\
Optimizer step (an event changes state) & event listeners (\lstinline{change}, \lstinline{drop}, \lstinline{click}) \\
Logging / metrics readout & \lstinline{drawChart()} on the canvas \\
\bottomrule
\end{tabular}
\end{center}

There is no framework (no React, no Vue). Every screen update is: something
happens $\to$ a small handler function mutates the \lstinline{state} object
$\to$ a \lstinline{render*()} function throws away the old HTML for that
section and rebuilds it from scratch. This is deliberately unsophisticated --
closer to re-running a plotting cell than to a real UI framework's diffing
engine -- and that is exactly what makes it easy to read top to bottom.

\section{Where Everything Lives}

\begin{center}
\begin{tabular}{@{}p{0.32\textwidth}p{0.6\textwidth}@{}}
\toprule
\textbf{File} & \textbf{What it's responsible for} \\
\midrule
\lstinline{playground/server.py} &
The backend. Stdlib-only Python (no Flask/FastAPI). Reads the CSV in
\lstinline{data/}, turns whatever the browser sends into a valid
\lstinline{config.yaml}, launches \lstinline{entrypoint.py} as a subprocess,
and parses its stdout into loss points. \\
\addlinespace
\lstinline{playground/static/index.html} &
The skeleton: every panel, slot, button, and input that exists, and nothing
else. Think of it as declaring the layout only -- like defining
\lstinline{nn.Module} submodules in \lstinline{__init__} before any
\lstinline{forward} logic exists. \\
\addlinespace
\lstinline{playground/static/app.js} &
All behavior. The \lstinline{state} object, every
\lstinline{render*()} function, all drag-and-drop wiring, the Run button,
polling, and the chart. \textbf{You will spend 90\% of your time in this
file.} \\
\addlinespace
\lstinline{playground/static/style.css} &
Pure appearance: colors, spacing, borders. Nothing in here changes what the
app \emph{does}. Safe to ignore unless something looks visually wrong. \\
\bottomrule
\end{tabular}
\end{center}

Start the whole thing with:
\begin{lstlisting}
python playground/server.py
\end{lstlisting}
It opens \lstinline{http://127.0.0.1:8765} in your browser automatically.

\section{The State Object: Your Single Source of Truth}

At the top of \lstinline{app.js} there is one object that fully describes
everything you have configured. Nothing else needs to be inspected to know
what the app currently thinks the pipeline is.

\begin{lstlisting}
const state = {
  csv: { filename: "", columns: [] },   // from GET /api/csv
  target: null,                         // e.g. "species"
  excludedFeatures: new Set(),          // columns you unchecked
  encoder: null,                        // { type, embedding_dim, hidden_dims, dropout }
  head: null,                           // { type, hidden_dims, dropout }
  training: {
    optimizer, loss_fn, splits, epochs,
    batch_size, learning_rate, device,
  },
};
\end{lstlisting}

Every control on screen either \emph{writes into} this object (an input's
\lstinline{change} handler) or \emph{reads from} it (a
\lstinline{render*()} function). If a number on screen is wrong, the bug is
in exactly one of those two directions -- either the wrong handler is writing
the wrong field, or the wrong field is being read when rendering.

\section{Screen-by-Screen Tour}

\subsection{Top Bar: Component Palette}
Two draggable squares: \emph{MLP Encoder} and \emph{Classification Head}.
These are the only two component types wired up on the Python side today
(see Section~\ref{sec:new-component} to add a third). Each palette block
carries two HTML \lstinline{data-*} attributes -- \lstinline{data-role}
(\lstinline{encoder} or \lstinline{head}) and \lstinline{data-type}
(\lstinline{mlp} or \lstinline{classification}) -- that get read in
\lstinline{wirePalette()} and written into the drag payload.

\subsection{Left Panel: CSV Columns and Feature Selection}
On load, \lstinline{init()} calls \lstinline{GET /api/csv}, which the
backend answers by looking at \lstinline{data/*.csv} (whichever CSV file it
finds first) and returning its header row. \lstinline{renderColumnStack()}
turns every column that isn't currently the target into a chip with:
\begin{itemize}
  \item a checkbox -- checked means ``use this column as a model input.''
  Unchecking it adds the column name to \lstinline{state.excludedFeatures}
  and greys out the chip (\lstinline{.excluded} CSS class) without removing
  it, so you can toggle it back on;
  \item a draggable label -- drag it onto the Target slot to make it the
  prediction target instead.
\end{itemize}
The feature list actually sent to the backend is computed fresh in
\lstinline{onRun()}: every CSV column, minus the target, minus anything in
\lstinline{excludedFeatures}.

\subsection{Center: The Playground Slots}
Three fixed slots, left to right: Encoder $\to$ Head $\to$ Target. Each slot
only accepts drops whose \lstinline{data-role} (or, for the target slot, drag
``kind'') matches its own role -- dropping a Head onto the Encoder slot is a
silent no-op with a one-line status message, not a crash.

Once a slot is filled, its \lstinline{render*Slot()} function
(\lstinline{renderEncoderSlot}, \lstinline{renderHeadSlot},
\lstinline{renderTargetSlot}) replaces the placeholder text with the
component's name, a $\times$ button to remove it, and -- for
Encoder/Head -- plain number/text inputs bound directly to that component's
hyperparameters. There is no ``settings panel'' or modal to open; the inputs
are always visible once the slot is filled.

\subsection{Right Panel: The Loss Chart}
A bare \lstinline{<canvas>} element. \lstinline{drawChart()} clears it and
redraws two polylines (train loss in blue, val loss in red) every time new
points arrive from polling. There is no charting library -- it is manual
pixel math (\lstinline{xScale}/\lstinline{yScale} map epoch/loss values to
pixel coordinates). If you want gridlines, axis labels, or a legend beyond
the two colored dots under the canvas, this is the function to extend.

\subsection{Bottom Bar: Training Controls}
Optimizer, loss function, train/val/test split, epochs, batch size, learning
rate, device, and the Run button. Every input has a plain
\lstinline{change} listener in \lstinline{wireBottomBar()} that writes
straight into \lstinline{state.training}. The split fields are \emph{not}
required to sum to 1 -- the backend normalizes them
(\lstinline{split / sum(splits)}) before writing \lstinline{config.yaml}.

\section{What Does the ``Hidden Dims'' Field Actually Accept?}
\label{sec:hidden-dims}

Both the Encoder and Head slots expose a plain text box for
\lstinline{hidden_dims}, not a structured widget. It accepts a
\textbf{comma-separated list of positive integers}, e.g.:

\begin{lstlisting}
32          ->  one hidden layer of width 32
128, 64     ->  two hidden layers, widths 128 then 64
(empty)     ->  no hidden layers at all (a single Linear layer)
\end{lstlisting}

The parsing happens in \lstinline{parseDims()} in \lstinline{app.js}: split
on commas, trim whitespace, drop empty pieces, convert to numbers, then
\textbf{silently discard anything that isn't a positive integer}. That last
filter was added deliberately: without it, typing something like
\lstinline{12,,7} or \lstinline{abc} would produce a \lstinline{NaN} in the
array. \lstinline{JSON.stringify} turns \lstinline{NaN} into \lstinline{null}
on the wire, which becomes \lstinline{None} in Python, which reaches
\lstinline{ClassificationHead}/\lstinline{TabularMLPEncoder} as a hidden
dimension and throws a \lstinline{TypeError} several stack frames away from
the actual typo. Now a bad entry is just quietly dropped instead.

Empty string $\to$ empty list \lstinline{[]} is intentional and safe: it
degenerates to a single \lstinline{nn.Linear(input_dim, output_dim)}, which
is exactly what \lstinline{ClassificationHead} and \lstinline{TabularMLPEncoder}
do internally when \lstinline{hidden_dims} is empty.

\section{What Happens When You Click Run}

This is the one flow worth memorizing, because it is the entire bridge
between the browser and your training code:

\begin{enumerate}
  \item \textbf{Validate (JS).} \lstinline{onRun()} checks an encoder, a
  head, and a target are all set, and that at least one feature column
  survives the checkboxes.
  \item \textbf{Serialize (JS).} The relevant slice of \lstinline{state} is
  packed into a plain JSON object and POSTed to \lstinline{/api/run}.
  \item \textbf{Build the config (Python).} \lstinline{build_config()} in
  \lstinline{server.py} fills in everything the UI doesn't expose
  (checkpointing defaults, computed \lstinline{input_dim}/\lstinline{output_dim},
  normalized splits) and returns a dict shaped exactly like
  \lstinline{DeepScratchConfig}.
  \item \textbf{Write YAML.} That dict is dumped straight over
  \lstinline{config.yaml} at the repo root -- the same file
  \lstinline{entrypoint.py} always reads.
  \item \textbf{Launch training.}
  \lstinline{subprocess.Popen(["python", "entrypoint.py"], cwd=REPO_ROOT)} --
  literally the same command you'd type by hand.
  \item \textbf{Stream stdout.} A background thread reads subprocess output
  one line at a time. \lstinline{BasicTorchTrainer} already prints one line
  per epoch (\lstinline{Epoch 3/50 - train_loss: 0.41}); a
  regex (\lstinline{EPOCH_LINE}) turns each matching line into a
  \lstinline|{epoch, train_loss, val_loss}| point. No changes to your
  training code were needed for this to work.
  \item \textbf{Poll.} The browser calls \lstinline{GET /api/status} every
  600ms and redraws the chart with whatever new points showed up.
\end{enumerate}

\section{Web-Dev-to-Python Glossary}

\begin{center}
\begin{tabular}{@{}p{0.22\textwidth}p{0.7\textwidth}@{}}
\toprule
\textbf{Term} & \textbf{What it means to you} \\
\midrule
DOM & The live tree of on-screen elements. Think of it as the object graph a
\lstinline{render*()} function builds and hands to the browser to paint. \\
Event listener & A registered callback, same idea as a hook -- ``when X
happens, call this function.'' \lstinline{addEventListener("change", fn)}
is not conceptually different from registering a backward hook on a tensor. \\
\lstinline{fetch} & Python's \lstinline{requests.get}/\lstinline{requests.post},
just built into the browser. \\
JSON payload & The dict you \lstinline{json.dumps()} and ship over HTTP --
identical in spirit to any REST API body you've already sent from Python. \\
Canvas & A blank pixel surface you draw on yourself with explicit coordinates,
the rough web equivalent of calling \lstinline{plt.plot} but drawing every
line segment manually. \\
\lstinline{dataTransfer} & The small string payload attached to a drag
gesture -- like passing a serialized argument into a callback that fires on
drop. \\
\lstinline{data-*} attribute & A free-form key/value tag on an HTML element,
used here purely to stash metadata (which component type a palette block
represents) for JS to read back later. \\
\bottomrule
\end{tabular}
\end{center}

\section{Cookbook: ``If You Want To \ldots''}

\begin{itemize}
  \item \textbf{Change a default hyperparameter value} (e.g. default
  embedding dim) $\to$ edit the \lstinline{DEFAULTS} object at the top of
  \lstinline{app.js}.
  \item \textbf{Add an optimizer/loss option to the dropdown} $\to$ add an
  \lstinline{<option>} in \lstinline{index.html}, \emph{and} make sure a
  matching key exists in the \lstinline{OPTIMIZERS}/\lstinline{LOSS_FNS}
  dict in \lstinline{entrypoint.py} -- the dropdown value is written verbatim
  into \lstinline{config.yaml} and must resolve to something real.
  \item \textbf{Change chart colors} $\to$ the hex strings passed into
  \lstinline{drawLine(ctx, key, color, ...)} calls inside
  \lstinline{drawChart()} in \lstinline{app.js} (and, for the legend dots
  under the chart, the matching colors in \lstinline{style.css}).
  \item \textbf{Change the polling rate} $\to$ the \lstinline{600} (ms)
  literal in \lstinline{poll()} in \lstinline{app.js}.
  \item \textbf{Change the server port} $\to$
  \lstinline{python playground/server.py <port>}, or edit the default in
  \lstinline{main()} in \lstinline{server.py}.
  \item \textbf{Expose a currently-hidden config field in the UI} (e.g.
  checkpointing options) $\to$ add a control to the bottom bar in
  \lstinline{index.html}, wire it in \lstinline{wireBottomBar()}, add it to
  \lstinline{state.training}, and read it in \lstinline{build_config()} in
  \lstinline{server.py} instead of the hardcoded default.
\end{itemize}

\section{Adding a Brand-New Component Type End to End}
\label{sec:new-component}

Say you implement a \lstinline{RegressionHead} in
\lstinline{deepscratch/heads/}. To make it appear in the playground:

\begin{enumerate}
  \item \textbf{Register it where training already knows about types.} Add a
  branch to \lstinline{build_head()} in \lstinline{entrypoint.py} for
  \lstinline{type == "regression"}. This step alone would already let you
  hand-edit \lstinline{config.yaml} and run \lstinline{entrypoint.py}
  directly -- the playground is just a front end onto this registry.
  \item \textbf{Add a palette block.} In \lstinline{index.html}:
  \begin{lstlisting}
<div class="block palette-block" draggable="true"
     data-role="head" data-type="regression">Regression Head</div>
  \end{lstlisting}
  \item \textbf{Give it defaults.} In \lstinline{app.js}, add an entry to
  \lstinline{DEFAULTS}:
  \begin{lstlisting}
regression: { hidden_dims: "", dropout: 0 },
  \end{lstlisting}
  \item \textbf{Render its controls.} If it needs the same
  hidden-dims/dropout inputs as \lstinline{ClassificationHead}, no extra
  work is needed -- \lstinline{renderHeadSlot()} is generic. If it needs a
  different input (e.g. no dropout), branch on \lstinline{state.head.type}
  inside \lstinline{renderHeadSlot()}.
  \item \textbf{Update} \lstinline{build_config()} \textbf{in
  \lstinline{server.py}} if the new type needs different fields written into
  \lstinline{components.head} (e.g. \lstinline{output_dim} should probably
  become \lstinline{1} for regression instead of the class count).
\end{enumerate}

The pattern is the same for a new encoder, trainer, or data module -- the
Python side (\lstinline{entrypoint.py}'s \lstinline{build_*} functions) is
the source of truth for what \emph{exists}; the playground only needs to
learn how to draw a box for it and what defaults to prefill.

\section{Quick-Reference Table}

\begin{center}
\begin{tabular}{@{}p{0.3\textwidth}p{0.28\textwidth}p{0.32\textwidth}@{}}
\toprule
\textbf{On-screen control} & \textbf{HTML id} & \textbf{Lives in / driven by} \\
\midrule
Encoder slot & \lstinline{encoder-slot} & \lstinline{renderEncoderSlot()} \\
Head slot & \lstinline{head-slot} & \lstinline{renderHeadSlot()} \\
Target slot & \lstinline{target-slot} & \lstinline{renderTargetSlot()} \\
Column stack & \lstinline{column-stack} & \lstinline{renderColumnStack()} \\
Loss chart & \lstinline{loss-chart} & \lstinline{drawChart()} \\
Status line & \lstinline{status-text} & \lstinline{setStatus()} \\
Run button & \lstinline{run-button} & \lstinline{onRun()} \\
Split inputs & \lstinline{split-train/val/test} & \lstinline{state.training.splits} \\
\bottomrule
\end{tabular}
\end{center}

\end{document}
